\chapter{O sistema dual: regra e reconstrução}
\label{cap:dual}

\begin{objetivos}
Ao final deste capítulo, você será capaz de:
\begin{itemize}[leftmargin=*,itemsep=0.2em]
  \item descrever a arquitetura de um sistema dual que combina um diagrama de decisão
        determinístico com um autoencoder recorrente;
  \item justificar por que a combinação supera cada componente isolado, em termos de
        cobertura e de taxa de falsos positivos;
  \item descrever as decisões de projeto de um autoencoder LSTM para detecção de
        anomalias e a definição de seu limiar pela regra $\mu + 3\sigma$;
  \item explicar como o sistema se autotreina em campo, sem rotulação humana, e quais
        salvaguardas isso exige;
  \item interpretar os resultados de três estudos de caso reais e compará-los com a
        literatura sobre fechamento espúrio de DHSV.
\end{itemize}
\end{objetivos}

\section{O problema que este capítulo resolve}

Comecemos pelo requisito, não pelo método.

Uma operadora deseja monitorar duzentos e cinquenta poços submarinos em tempo real.
Cada poço tem um comportamento distinto: níveis de pressão diferentes, curvas de
declínio diferentes, esquemas de elevação artificial diferentes. Não há rotulação
humana disponível de forma contínua. Não há garantia de que os eventos que ocorrerão
estejam catalogados. E o sistema precisa funcionar a partir do momento em que é ligado.

Essas quatro restrições eliminam, de saída, a solução acadêmica padrão --- treinar um
classificador supervisionado sobre um conjunto rotulado e implantá-lo. O que se propõe
neste capítulo, e que efetivamente opera em campo, é uma arquitetura de dois
componentes com naturezas complementares.

\begin{description}[leftmargin=0pt,style=nextline,itemsep=0.5em]

\item[Um diagrama de decisão determinístico.] Codifica conhecimento de engenharia em
regras explícitas. Não requer treinamento, funciona no primeiro segundo, é auditável e
seu comportamento é previsível. Em contrapartida, só reconhece o que foi explicitamente
programado.

\item[Um autoencoder recorrente.] Aprende o comportamento normal \emph{daquele poço
específico} a partir dos próprios dados, sem rótulos, e sinaliza qualquer desvio.
Detecta o inesperado, mas não sabe nomeá-lo e precisa de tempo para aprender.

\end{description}

\begin{destaque}[A tese deste capítulo]
Os dois componentes não competem: eles se alimentam. O diagrama de decisão determina
\emph{quais dados são normais} e, com isso, fornece ao autoencoder o conjunto de
treinamento que ele não teria de outra forma. O autoencoder, por sua vez, cobre o
espaço de eventos que o diagrama não enumera. É uma relação de mão dupla, e é essa
relação que produz o ganho de desempenho reportado na \cref{sec:dual-resultados}.
\end{destaque}

\section{Arquitetura geral}

\begin{figure}[htbp]
\centering
\includegraphics[width=0.55\textwidth]{DiagramaVIGIA2.png}
\caption{Fluxo geral do sistema dual. Os dados do poço alimentam simultaneamente o
diagrama de decisão e o módulo de detecção de anomalias por inteligência artificial. O
diagrama de decisão encaminha ao módulo de aprendizado apenas os trechos que classifica
como normais. Qualquer um dos dois componentes pode disparar o alerta ao operador.
Adaptado de \citet{aranha2023system}.}
\label{fig:dual-visaogeral}
\end{figure}

A \cref{fig:dual-visaogeral} mostra a estrutura. Note a seta rotulada como entrada de
dados normais: é ela que materializa a relação descrita acima. O autoencoder nunca vê
dados que o diagrama de decisão considera anômalos, e portanto o modelo de normalidade
que ele constrói não é contaminado por eventos.

\section{O diagrama de decisão}

O diagrama de decisão é uma árvore de regras organizada em quatro camadas sucessivas.
Cada camada responde a uma pergunta e restringe o espaço de hipóteses da seguinte.

\begin{figure}[htbp]
\centering
\includegraphics[width=0.9\textwidth]{New_Decision_Tree.jpeg}
\caption{Estrutura do diagrama de decisão. O fluxo bifurca conforme o dado seja o
primeiro do poço, siga uma mudança de estado de válvula ou dê continuidade a um regime já
estabelecido, encaminhando cada caso para estabilização, análise de anomalia ou
atualização da média móvel. Adaptado de \citet{aranha2023system}.}
\label{fig:dual-decisao}
\end{figure}

A \cref{fig:dual-decisao} mostra o fluxo. Vale insistir num ponto que a semelhança visual
esconde: apesar do nome e do formato, isto \emph{não} é uma árvore de decisão aprendida.
Nenhum de seus nós foi ajustado por otimização; cada bifurcação foi escrita por
engenheiros de produção e pode ser lida, auditada e alterada linha a linha. É conhecimento
declarado, não induzido --- e é essa diferença que sustenta a arquitetura dual descrita
neste capítulo.

\begin{enumerate}[leftmargin=*,itemsep=0.5em]

\item \textbf{Identificação do poço.} Qual poço é este, qual seu esquema de completação,
quais sensores estão disponíveis e quais são os limites operacionais declarados? Essa
camada não classifica nada --- ela carrega o contexto sem o qual as camadas seguintes
não podem operar. É também onde entram informações que nenhum modelo aprende dos dados:
profundidade, número de zonas produtoras, presença de gás de elevação.

\item \textbf{Tipo de dado.} O sinal recebido é válido? Esta camada implementa as
verificações discutidas no \cref{cap:dominio}: valor ausente, sensor congelado, leitura
fora da faixa física do instrumento, salto incompatível com a dinâmica do processo. Um
dado reprovado aqui não prossegue, e não contamina nem o alerta nem o treinamento.

\item \textbf{Transiente ou estabilizado?} O poço está em regime permanente ou em
transição? A distinção é essencial: durante uma partida, uma parada programada ou uma
manobra de válvula, os sinais apresentam variações de grande amplitude que são
perfeitamente normais. Sinalizá-las como anomalia é a receita para a fadiga de alarme
descrita no \cref{cap:introducao}.

\item \textbf{Análise de anomalia.} Somente aqui, com poço identificado, dado válido e
regime estabilizado, aplicam-se as regras de diagnóstico propriamente ditas ---
divergência entre pressões, queda de temperatura, oscilação sustentada.

\end{enumerate}

\begin{destaque}[Por que quatro camadas e não uma lista de regras]
A organização em camadas não é estética. Cada camada elimina uma fonte de falso
positivo antes que a próxima seja avaliada, e o resultado é que a última camada ---
a única que gera alarme --- opera sobre um subconjunto de dados já saneado. Uma lista
plana de regras avaliadas em paralelo produziria alarmes durante partidas, durante
falhas de sensor e durante manobras programadas.
\end{destaque}

\subsection{O tratamento de manobras de válvula}

Uma consequência prática merece destaque, porque é onde a maior parte dos sistemas
falha. Quando o operador altera deliberadamente o estado de uma válvula, o poço muda de
ponto de operação. Todos os sinais se deslocam. Um detector de anomalia baseado em
desvio da normalidade aprendida sinalizará --- corretamente, do ponto de vista
estatístico, e inutilmente, do ponto de vista operacional.

A solução adotada é radical e correta: \emph{ao detectar mudança de válvula, o sistema
descarta o modelo aprendido e reinicia o ciclo de treinamento}. O poço, depois da
manobra, é tratado como um poço novo. O custo é o período de reaprendizado; o benefício
é a ausência de alarmes espúrios sistemáticos.

\section{O detector por autoencoder}

\subsection{Lógica de operação}

\begin{figure}[htbp]
\centering
\includegraphics[width=\textwidth]{IA_flow_chart_2.jpeg}
\caption{Lógica operacional do módulo de detecção por inteligência artificial. Note os
três estados possíveis: antes do treinamento, o sistema apenas acumula dados normais
selecionados pelo diagrama de decisão; após o treinamento, ele prediz, compara com a
tolerância e informa o estado do poço; uma mudança de válvula reinicia o ciclo. Dados
julgados anômalos nunca retornam ao conjunto de treinamento. Adaptado de
\citet{aranha2023system}.}
\label{fig:dual-fluxo-ia}
\end{figure}

A \cref{fig:dual-fluxo-ia} detalha a máquina de estados. Três aspectos merecem
comentário.

\begin{itemize}[leftmargin=*,itemsep=0.4em]

\item \textbf{O sistema informa que não sabe.} Antes de acumular dados suficientes, ele
responde explicitamente ``rede não treinada''. Não há resposta silenciosamente errada.
É uma forma primitiva --- mas efetiva --- da rejeição que o \cref{cap:binarios}
formalizará.

\item \textbf{Dados anômalos são excluídos do treinamento futuro.} Sem essa
salvaguarda, o modelo aprenderia gradualmente a considerar a anomalia normal, e o
alarme se extinguiria sozinho. É a versão local do problema de deriva discutido no
\cref{cap:operacao}.

\item \textbf{O treinamento é contínuo.} Dados normais posteriores ao treinamento
inicial continuam sendo acumulados e usados em retreinamentos, o que permite ao modelo
acompanhar a evolução lenta e legítima do poço.

\end{itemize}

\subsection{Arquitetura da rede}

O autoencoder é do tipo \ing{sequence-to-sequence} com camadas LSTM. A entrada é um
tensor de forma $(\text{lote}, 500, 5)$: quinhentos passos de tempo de cinco sensores.

\begin{center}
\small
\begin{tabular}{lll}
\toprule
\textbf{Camada} & \textbf{Forma de saída} & \textbf{Função} \\
\midrule
Entrada                       & $(500, 5)$  & janela de \SI{500}{\second} \\
LSTM(16), ReLU                & $(500, 16)$ & codificação \\
LSTM(4), ReLU                 & $(4)$       & \textbf{representação latente} \\
\ing{RepeatVector}(500)       & $(500, 4)$  & replicação temporal \\
LSTM(16), ReLU                & $(500, 16)$ & decodificação \\
\ing{TimeDistributed} Densa(5)& $(500, 5)$  & reconstrução \\
\bottomrule
\end{tabular}
\end{center}

A compressão é de $500 \times 5 = \num{2500}$ valores para $4$: um fator de $625$. Essa
severidade é deliberada, e responde diretamente à advertência do \cref{cap:profundo}
sobre autoencoders excessivamente expressivos. Com quatro dimensões latentes, a rede não
tem como memorizar; ela é obrigada a capturar apenas a estrutura dominante da
normalidade.

\subsection{Treinamento}

\begin{center}
\small
\begin{tabular}{ll}
\toprule
\textbf{Hiperparâmetro} & \textbf{Valor} \\
\midrule
Função de perda      & SMAPE (\cref{eq:smape}) \\
Otimizador           & Adam \\
Épocas               & 500 \\
Tamanho do lote      & 100 \\
Ativação             & ReLU \\
Normalização         & Min-Max para $[0,1]$ \\
Ruído injetado       & 0 a 0,5\,\% \\
Dados mínimos        & $\approx 500$ amostras normais ($\approx \SI{5000}{\second}$) \\
\bottomrule
\end{tabular}
\end{center}

Duas escolhas merecem justificativa.

A \textbf{perda SMAPE} em vez do erro quadrático médio decorre da heterogeneidade de
escalas discutida no \cref{cap:aprendizado}: pressões da ordem de centenas de bar e
temperaturas da ordem de dezenas de graus. Um erro quadrático médio sobre variáveis
normalizadas ainda pondera implicitamente pela variância de cada canal; o SMAPE, sendo
relativo, trata todos os canais de forma comparável.

O \textbf{ruído injetado} de até 0,5\,\% cumpre dois papéis. Regulariza, impedindo que a
rede se apoie em detalhes de alta frequência específicos daquele poço; e simula a
variabilidade de instrumentação, tornando o modelo mais robusto quando transferido
entre poços com sensores de calibração ligeiramente diferente.

\subsection{Limiar e classificação de risco}

O escore de anomalia é o SMAPE de reconstrução. Sobre o conjunto de dados normais
acumulados, calculam-se $\mu$ e $\sigma$ desse escore, e define-se
\begin{equation}
  \tau = \mu + 3\sigma.
  \label{eq:dual-limiar}
\end{equation}

Sob normalidade aproximada, $\tau$ corresponde a aproximadamente $99{,}7\,\%$ dos casos
--- uma taxa nominal de falso positivo de $0{,}3\,\%$. A distribuição do erro de
reconstrução não é gaussiana, mas a regra é adequada como heurística e tem a virtude de
ser explicável a um operador em uma frase.

O valor do escore é então convertido em três níveis de risco:

\begin{figure}[htbp]
\centering
\includegraphics[width=0.86\textwidth]{abnormal.png}
\caption{Exemplo de detecção em dados reais. Nos dois painéis superiores, as pressões e
as temperaturas medidas por PDG e TPT; no inferior, o escore de anomalia com a linha de
tolerância da \cref{eq:dual-limiar}. A ultrapassagem do limiar precede em várias horas o
desvio que seria visível a olho nu nos sinais brutos. Adaptado de
\citet{aranha2023system}.}
\label{fig:dual-escore}
\end{figure}

\begin{center}
\small
\begin{tabular}{ll l}
\toprule
\textbf{Faixa} & \textbf{Nível} & \textbf{Ação recomendada} \\
\midrule
$< 30\,\%$          & Baixo & Registro; sem notificação. \\
$30\,\%$ a $70\,\%$ & Médio & Notificação ao operador; acompanhamento. \\
$> 70\,\%$          & Alto  & Alerta; investigação imediata. \\
\bottomrule
\end{tabular}
\end{center}

\begin{destaque}[Por que graduar o alarme]
Um alarme binário força uma escolha entre perder eventos e cansar o operador. A
graduação em três níveis desloca essa escolha para o operador, que dispõe de contexto
que o sistema não tem. É uma decisão de projeto de interface, não de aprendizado de
máquina, e é frequentemente o que separa um sistema usado de um sistema desligado.
\end{destaque}

\section{Implantação}

\begin{figure}[htbp]
\centering
\includegraphics[width=\textwidth]{system.png}
\caption{Interface do sistema em operação. À esquerda, os valores instantâneos dos
cinco sensores monitorados; ao centro, os registros de pressão e temperatura; à
direita, os erros relativos de reconstrução por sensor e o erro agregado. O painel
inferior direito mostra o escore de anomalia cruzando o limiar, indicado pela linha
horizontal. Adaptado de \citet{aranha2023system}.}
\label{fig:dual-interface}
\end{figure}

O sistema foi implementado com Node-RED \citep{nodered} para orquestração de fluxo e
aquisição, e Python 3.7 com Keras \citep{Chollet2015} para o modelo. A escolha do
Node-RED merece nota: trata-se de uma ferramenta de programação visual por fluxo,
amplamente usada em automação industrial, que permite a engenheiros de operação
inspecionar e modificar o encadeamento sem tocar no código do modelo.

À época da avaliação reportada aqui, o sistema estava implantado em \textbf{mais de
vinte unidades flutuantes de produção} e \textbf{mais de duzentos e cinquenta poços
submarinos} nas bacias de Santos e de Campos.

\begin{destaque}[O que ``implantado'' significa]
Vale a pena explicitar, porque a palavra é usada com generosidade na literatura. Aqui
significa: executando continuamente em ambiente de produção da operadora, sobre dados
de tempo real, com saída visível a operadores em turno, e sujeito às consequências
operacionais de seus alertas. Não é um protótipo avaliado \ing{offline}.
\end{destaque}

\section{Resultados}
\label{sec:dual-resultados}

\subsection{Contribuição de cada componente}

A \cref{tab:dual-componentes} isola o efeito de acrescentar o autoencoder ao diagrama de
decisão.

\begin{table}[htbp]
\centering
\caption{Desempenho do diagrama de decisão isolado e do sistema dual completo.}
\label{tab:dual-componentes}
\begin{tabular}{lccc}
\toprule
\textbf{Configuração} & \textbf{ACC} & \textbf{ACC}$_b$ & $\bm{F_1}$ \\
\midrule
Diagrama de decisão                 & $0{,}9727$ & $0{,}9615$ & $0{,}9758$ \\
Diagrama de decisão + autoencoder   & $\bm{0{,}9894}$ & $\bm{0{,}9771}$ & $\bm{0{,}9917}$ \\
\bottomrule
\end{tabular}
\end{table}

\begin{resultadochave}
O acréscimo do autoencoder reduz o erro de $2{,}73\,\%$ para $1{,}06\,\%$ --- uma
redução relativa de aproximadamente $61\,\%$. A acurácia balanceada sobe de $0{,}9615$
para $0{,}9771$, e a medida $F_1$ de $0{,}9758$ para $0{,}9917$.
\end{resultadochave}

A leitura correta desse ganho exige cuidado. Em valor absoluto, $1{,}67$ ponto
percentual parece pouco. Em termos operacionais, significa que aproximadamente três em
cada cinco erros do sistema baseado em regras deixaram de ocorrer. Quando a base é um
sistema já bom, é assim que o progresso se manifesta.

\subsection{Comparação com a literatura}

O fechamento espúrio de DHSV (classe 2) é o evento mais estudado na literatura do 3W,
o que permite comparação direta.

\begin{table}[htbp]
\centering
\caption{Detecção de fechamento espúrio da DHSV: comparação com trabalhos anteriores.}
\label{tab:dual-literatura}
\begin{tabular}{llcc}
\toprule
\textbf{Trabalho} & \textbf{Método} & \textbf{ACC} & $\bm{F_1}$ \\
\midrule
\citet{Marins2021}  & Floresta aleatória      & $0{,}8708$ & --- \\
\citet{Turan2021}   & Árvore de decisão       & $0{,}6000$ & $0{,}4900$ \\
\citet{Machado2022} & LSTM                    & $\bm{0{,}9922}$ & $0{,}9360$ \\
Este capítulo       & Diagrama + autoencoder  & $0{,}9894$ & $\bm{0{,}9917}$ \\
\bottomrule
\end{tabular}
\end{table}

\begin{destaque}[Como ler esta tabela]
O sistema dual tem acurácia ligeiramente inferior à de \citet{Machado2022}
($0{,}9894$ contra $0{,}9922$) e medida $F_1$ substancialmente superior
($0{,}9917$ contra $0{,}9360$).

A combinação é diagnóstica. Uma medida $F_1$ de $0{,}9360$ com acurácia de $0{,}9922$
indica desequilíbrio entre precisão e revocação --- tipicamente, muitos falsos
positivos ou muitos falsos negativos concentrados na classe minoritária, mascarados
pela acurácia global (\cref{cap:aprendizado}). O sistema dual apresenta as duas métricas
alinhadas, o que indica desempenho equilibrado.

Uma ressalva de honestidade: os conjuntos de avaliação não são idênticos entre os
trabalhos, e a comparação é indicativa, não conclusiva.
\end{destaque}

\subsection{Estudos de caso em poços reais}

Os três casos a seguir são eventos reais, com data e hora, detectados pelo sistema em
operação. Eles importam porque são a única evidência disponível sobre generalização a
condições que nenhum conjunto de dados público contém.

\begin{table}[htbp]
\centering
\caption{Estudos de caso: características dos poços. WD é a lâmina d'água e TD a
profundidade total.}
\label{tab:dual-pocos}
\small
\begin{tabular}{lllll}
\toprule
\textbf{Poço} & \textbf{Tipo} & \textbf{WD (m)} & \textbf{TD (m)} & \textbf{Zonas} \\
\midrule
1 & Produtor, pré-sal de Santos & \num{2041} & \num{5870} & 3 (\num{5432} a \num{5774}\,m) \\
2 & Produtor                    & \num{2125} & \num{5552} & 2 (\num{5137} a \num{5440}\,m) \\
3 & Injetor WAG                 & \num{1946} & \num{6197} & 3 (\num{5437} a \num{5730}\,m) \\
\bottomrule
\end{tabular}
\end{table}

\begin{table}[htbp]
\centering
\caption{Estudos de caso: eventos detectados e desempenho. O SMAPE indicado é o valor de
pico do erro de reconstrução durante o evento.}
\label{tab:dual-casos}
\small
\begin{tabular}{llcccc}
\toprule
\textbf{Poço} & \textbf{Evento} & \textbf{Data e hora} & \textbf{SMAPE} &
\textbf{ACC}$_b$ & $\bm{F_1}$ \\
\midrule
1 & Fechamento espúrio (1º) & 01/08/2022 01:16 & $39{,}148\,\%$ & $0{,}9897$ & $0{,}9987$ \\
1 & Fechamento espúrio (2º) & 01/08/2022 08:48 & $48{,}312\,\%$ & $0{,}9988$ & $0{,}9989$ \\
2 & Comunicação de pressão na CKGL & 06/06/2018 14:37 & $33{,}136\,\%$ & $0{,}9998$ & $0{,}9998$ \\
3 & Fechamento espúrio DHSV/M1 & 18/10/2022 02:10 & $23{,}390\,\%$ & $0{,}9988$ & $0{,}9991$ \\
\bottomrule
\end{tabular}
\end{table}

\begin{figure}[htbp]
\centering
\includegraphics[width=\textwidth]{fig8.png}
\caption{Poço 1, primeiro evento. De cima para baixo: os sensores PDG, PT e TPT, o estado
das válvulas e a comparação entre a classificação dos dois modelos (\emph{LSTM status} e
\emph{analytical status}) e a classificação humana posterior (\emph{class}). Adaptado de
\citet{aranha2023system}.}
\label{fig:dual-poco1a}
\end{figure}

\begin{figure}[htbp]
\centering
\includegraphics[width=\textwidth]{fig9.png}
\caption{Poço 1, segundo evento, no mesmo dia. Adaptado de \citet{aranha2023system}.}
\label{fig:dual-poco1b}
\end{figure}

A \cref{fig:dual-poco1a,fig:dual-poco1b} apresentam os dois eventos do poço 1, e a
comparação entre elas contém a lição mais útil do capítulo sobre a divisão de trabalho
entre as duas camadas.

No primeiro evento, às 01:16, \textbf{apenas o diagrama analítico detectou}. A rede não
teve chance: o fechamento espúrio da DHSV alterou o estado das válvulas, e o sistema
entrou em modo de análise --- justamente o comportamento projetado para evitar alarmes
durante manobras. O operador corrigiu com ciclagem de válvulas entre 01:39 e 01:52.

No segundo evento, às 08:48, \textbf{ambos detectaram}, e o diagrama classificou
corretamente a causa como fechamento espúrio da DHSV ou da M1.

\begin{destaque}[Por que isto justifica a arquitetura dual]
O primeiro evento é o argumento inteiro deste capítulo em uma linha: houve uma anomalia
real que a rede neural, por construção, não podia enxergar. A regra explícita enxergou.

Um sistema apenas aprendido teria perdido esse evento. Um sistema apenas analítico teria
perdido a classificação mais fina que a rede oferece em outros casos. A redundância não
é conservadorismo --- é a única configuração em que a janela cega de um componente é
coberta pelo outro.
\end{destaque}

\begin{figure}[htbp]
\centering
\includegraphics[width=\textwidth]{fig10.png}
\caption{Poço 2, retroanálise do evento de 6 de junho de 2018. Adaptado de
\citet{aranha2023system}.}
\label{fig:dual-poco2}
\end{figure}

\begin{figure}[htbp]
\centering
\includegraphics[width=\textwidth]{fig11.png}
\caption{Poço 3, evento de 18 de outubro de 2022, com o ciclo completo: detecção,
fechamento pelo operador, ciclagem, reabertura e retorno à normalidade. Adaptado de
\citet{aranha2023system}.}
\label{fig:dual-poco3}
\end{figure}

\begin{exemplo}{O poço 3, minuto a minuto}
O caso do poço 3, mostrado na \cref{fig:dual-poco3}, é o mais instrutivo, porque tem
começo, meio e fim documentados.

Trata-se de um injetor de água e gás alternados (WAG) em lâmina d'água de
\SI{1946}{\meter}. Às \textbf{02:10} de 18 de outubro de 2022, DHSV e M1 fecharam sem
comando. O sistema sinalizou o desvio com SMAPE de pico de $23{,}390\,\%$. As válvulas
foram reabertas às \textbf{07:14}, e o poço retornou à condição normal às
\textbf{08:29}.

Dois pontos merecem atenção. Primeiro, o SMAPE de pico deste caso é o menor dos quatro
--- cerca de metade do observado no segundo evento do poço 1 --- e ainda assim
suficiente para ultrapassar o limiar com folga. Isso indica que $\tau = \mu + 3\sigma$
está calibrado de forma conservadora o bastante para eventos de assinatura mais
discreta.

Segundo, o sistema acompanhou corretamente o \emph{retorno} à normalidade, com uma hora
e quinze minutos de defasagem em relação à reabertura das válvulas --- o tempo físico
de reestabilização do poço, não um atraso do algoritmo.
\end{exemplo}

\begin{exemplo}{O poço 2 e o valor de detectar o não catalogado}
O evento do poço 2 foi uma comunicação indevida de pressão através da válvula
reguladora de gás de elevação. \emph{Esse evento não pertence a nenhuma das nove classes
do 3W}, não estava programado no diagrama de decisão e não teria sido detectado por
nenhum classificador supervisionado treinado no catálogo padrão.

Foi detectado pelo autoencoder, com $\text{ACC}_b = 0{,}9998$, simplesmente porque não
se parecia com a normalidade daquele poço. A \cref{fig:dual-poco2} mostra a retroanálise:
o desvio é visível nos três sensores muito antes de a intervenção ter sido decidida em
campo.

É a demonstração mais direta do argumento deste livro: o componente que não sabe o nome
de nada é o único capaz de reagir ao que ainda não tem nome. O que lhe falta ---
dizer \emph{o que} está acontecendo --- é o que os \cref{cap:binarios,cap:mundoaberto}
e, em outro registro, a \cref{parte:explicabilidade} vão perseguir.
\end{exemplo}

\section{Limitações}

\begin{enumerate}[leftmargin=*,itemsep=0.4em]

\item \textbf{Não classifica.} O autoencoder responde ``anômalo'' ou ``normal''. O
diagrama de decisão classifica, mas apenas o que foi programado. Nenhum dos dois nomeia
um evento inédito.

\item \textbf{Período cego.} São necessárias aproximadamente quinhentas amostras normais
--- cerca de uma hora e vinte e três minutos --- antes que o modelo esteja disponível. E
o relógio reinicia a cada mudança de válvula.

\item \textbf{Limiar único e global.} Um único $\tau$ para todos os sensores e todos os
regimes. Um evento que afete fortemente um sensor e pouco os demais pode ser diluído no
erro agregado.

\item \textbf{Hipótese de normalidade sobre o erro.} A regra $\mu + 3\sigma$ pressupõe
uma distribuição que os dados não seguem exatamente. Na prática funciona, mas a taxa de
falso positivo nominal de $0{,}3\,\%$ não deve ser levada ao pé da letra.

\item \textbf{Avaliação limitada.} Três poços e quatro eventos. É evidência forte de
funcionamento e evidência fraca de generalização.

\end{enumerate}

\begin{destaque}[Para onde isso leva]
A limitação (1) é a mais estrutural, e é ela que abre o \cref{cap:binarios}: como
transformar ``algo está errado'' em ``isto é um hidrato'' \emph{sem} recair na hipótese
de mundo fechado que o \cref{cap:introducao} descartou.
\end{destaque}

\resumocapitulo{%
O sistema dual combina um diagrama de decisão determinístico em quatro camadas ---
identificação do poço, validação do dado, discriminação de transiente e análise de
anomalia --- com um autoencoder LSTM que comprime janelas de quinhentos segundos e cinco
sensores em quatro dimensões latentes. A relação entre os componentes é de mão dupla: o
diagrama seleciona os dados normais que treinam o autoencoder, e o autoencoder cobre os
eventos que o diagrama não enumera. O limiar de detecção é $\mu + 3\sigma$ sobre o SMAPE
de reconstrução, convertido em três níveis de risco. O sistema opera em mais de vinte
unidades flutuantes e mais de duzentos e cinquenta poços. O acréscimo do autoencoder ao
diagrama eleva a acurácia de $0{,}9727$ para $0{,}9894$ e a medida $F_1$ de $0{,}9758$
para $0{,}9917$, resultado competitivo com o estado da arte na detecção de fechamento
espúrio de DHSV. Quatro eventos reais em três poços confirmam o funcionamento em campo,
inclusive para um evento fora do catálogo. A limitação decisiva é que o sistema detecta
sem classificar --- e é essa lacuna que o restante da Parte II ataca.}

\begin{exercicios}
\item O autoencoder comprime $500 \times 5 = \num{2500}$ valores em quatro dimensões.
      Calcule a taxa de compressão e discuta, em termos de teoria da informação, que
      tipos de estrutura podem sobreviver a essa compressão e quais necessariamente se
      perdem.

\item Implemente o autoencoder descrito neste capítulo e treine-o sobre instâncias
      normais do 3W. Compare o histograma do SMAPE de reconstrução em dados normais e
      anômalos e verifique empiricamente se $\mu + 3\sigma$ é uma escolha adequada de
      limiar. Teste também $\mu + 2\sigma$ e $\mu + 4\sigma$, reportando a curva ROC.

\item O sistema reinicia o treinamento a cada mudança de válvula. Modele o custo dessa
      política: dado um poço com $\lambda$ manobras por mês e um período cego de
      \SI{5000}{\second}, qual a fração de tempo em que o sistema está indisponível?
      Para que valor de $\lambda$ a indisponibilidade ultrapassa $10\,\%$?

\item Proponha e justifique uma alternativa ao reinício completo após manobra de
      válvula. Considere: um modelo por estado de válvula; um modelo condicionado ao
      estado como entrada adicional; adaptação incremental. Discuta prós e contras.

\item O sistema usa um limiar único sobre o erro agregado. Proponha um esquema de
      limiares por sensor e discuta como combinar as decisões individuais. Que ganho
      você esperaria para um evento que afeta apenas a temperatura de fundo?

\item A \cref{tab:dual-literatura} mostra $\text{ACC} = 0{,}9922$ com $F_1 = 0{,}9360$
      para um dos trabalhos. Construa uma matriz de confusão compatível com esse par de
      valores e explique qual desequilíbrio ele revela.

\item Discuta criticamente a evidência dos estudos de caso. Quatro eventos em três poços
      sustentam qual afirmação? Formule três afirmações de força decrescente e indique
      qual delas os dados suportam.
\end{exercicios}

\begin{leituras}
\item \citet{LopesTese2025} --- a tese que originou este capítulo, com detalhamento
      adicional do diagrama de decisão.
\item \citet{Machado2022}, \citet{Marins2021} e \citet{Turan2021} --- os trabalhos
      comparados na \cref{tab:dual-literatura}.
\item \citet{aranha2023system} e \citet{Aranha2024} --- descrição de sistemas de
      monitoramento em operação, com os requisitos de campo que motivaram a arquitetura
      dual.
\item \citet{Malhotra2015LongST} --- referência clássica sobre autoencoders LSTM para
      detecção de anomalias em séries temporais.
\end{leituras}
